% Generated draft: requires mathematical and source review.
\documentclass[11pt]{article}
\usepackage[margin=25mm]{geometry}
\usepackage{fontspec}
\setmainfont{Latin Modern Roman}
\usepackage{amsmath,amssymb,mathtools,unicode-math}
\setmathfont{Latin Modern Math}
\usepackage{enumitem,needspace,xurl,hyperref}
\setlength{\parindent}{0pt}
\setlength{\parskip}{5pt}
\setlength{\emergencystretch}{3em}
\newcommand{\N}{\mathbb{N}}
\newcommand{\Z}{\mathbb{Z}}
\newcommand{\Q}{\mathbb{Q}}
\newcommand{\R}{\mathbb{R}}
\newcommand{\C}{\mathbb{C}}
\newcommand{\F}{\mathbb{F}}
\newcommand{\A}{\mathbb{A}}
\newcommand{\PP}{\mathbb{P}}
\newcommand{\E}{\mathbb{E}}
\newcommand{\eps}{\varepsilon}
\newcommand{\dd}{\,\mathrm{d}}
\begin{document}
\section*{UnsolvedMath ID: 30001880}
\textbf{Source problem code:} catalogue entry
\subsection{Definitions and mathematical statement}
Let \(Q=[0,1]^2\subset\R^2\), and let \(\lambda_2\) denote planar Lebesgue measure. A finite point set \(P\subseteq Q\) meets a set \(K\) if \(P\cap K\ne\varnothing\).

\textbf{Conditional formulation.} The following formulation assumes that the compact convex sets in the supplied statement are required to be contained in the unit square. This containment hypothesis is not explicit in the supplied wording. For \(0<\varepsilon\le1\), put
\[
\mathcal C_\varepsilon=\{K\subseteq Q:K\text{ is compact and convex, and }\lambda_2(K)\ge\varepsilon\},
\]
and define, whenever the minimum exists,
\[
f(\varepsilon)=\min\bigl\{|P|:P\subseteq Q\text{ is finite and }P\cap K\ne\varnothing\text{ for every }K\in\mathcal C_\varepsilon\bigr\}.
\]
The question is whether \(f(\varepsilon)=O(\varepsilon^{-1})\) as \(\varepsilon\downarrow0\). Explicitly, does there exist \(C>0\) and \(\varepsilon_0\in(0,1]\) such that for every \(\varepsilon\in(0,\varepsilon_0]\) there exists a finite \(P_\varepsilon\subseteq Q\) satisfying
\[
|P_\varepsilon|\le\frac{C}{\varepsilon}
\quad\text{and}\quad
\forall K\in\mathcal C_\varepsilon,\qquad P_\varepsilon\cap K\ne\varnothing?
\]
Here \(P_\varepsilon\) may depend on \(\varepsilon\), but must meet every member of \(\mathcal C_\varepsilon\) simultaneously.

\textbf{Unresolved source qualifications.} The supplied statement concerns all compact convex sets, whereas the background and literature assessment describe an original rectangle piercing question. These are different target families, and the supplied evidence does not resolve which is intended. No rectangle formulation is substituted here. The closed-square boundary convention and the containment hypothesis above require confirmation against the source. Without containment or another restriction relating the sets to the square, the wording specifies a different family.
\subsection{Short English statement}
Under the explicitly conditional interpretation that the target sets are compact convex subsets of the closed unit square, ask whether at most \(C/\varepsilon\) points can simultaneously meet every target of area at least \(\varepsilon\), with \(C\) independent of sufficiently small positive \(\varepsilon\). The source leaves containment implicit and conflicts with its rectangle-specific literature discussion, so the intended formulation requires review.
\subsection{Sources}
\begin{itemize}
\item Local UnsolvedMath registry: supplied complete record with permanent ID \(30001880\), source problem code \texttt{catalogue entry}, title ``Piercing Large Convex Sets in the Unit Square,'' in the Oberwolfach Reports Open Problems collection. The supplied proposed year is 2011; no individual proposer is supplied.
\item Supplied source citation: \emph{Discrete Geometry (2012)}, Oberwolfach Reports, DOI 10.4171/owr/2011/44, pp.~2459--2548. Original OWR report: \url{https://doi.org/10.4171/owr/2011/44}. The background describes these reports as records of mathematical workshops and their open-problem discussions. The citation's 2012 label and its 2011 DOI/proposed-year information are retained as supplied.
\item Supplied literature source: \emph{Piercing quasi-rectangles}, \url{https://math.nyu.edu/~pach/publications/quasirectangle031112.pdf}. The supplied assessment attributes the cited result to Pach--Tardos; complete bibliographic details are not supplied.
\item Literature assessment checked at \texttt{2026-08-21}: supplied dated curation, not independently verified. It labels the original rectangle piercing question open and states that the cited Pach--Tardos result gives a larger order for a relaxed quasi-rectangle class and therefore does not disprove the exact rectangle claim. The background expressly calls this dated literature triage derived from the supplied OWR curation and requests independent verification of time-sensitive status and cited claims. Its relation to the all-convex-set statement remains unresolved.
\item The registry labels the problem \texttt{open} and the statement \texttt{corrected\_verified}. Its verification note says only that TeX notation was enclosed in math delimiters without changing mathematical content; this does not establish mathematical or source verification.
\end{itemize}
\end{document}
